% Source: https://artofproblemsolving.com/community/c3h1726581_11_proved_with_functions
% Stable identity: aops-topic-1726581-q01
Every step of the argument is a correct deduction from the hypotheses: putting $x=0$ gives $1=n-f(-n)$, so $f(-n)=n-1$ for every integer $n$, and $n=0$ gives $f(0)=-1$. The conclusion $1=-1$ is false, so the \emph{hypotheses} are false: no function satisfies both conditions.

Indeed, $n=0$ in the defining equation already gives $f(x)=0-f(x)$, so $f(x)=0$ for every $x$; this contradicts $f(0)=1$. (Even the condition $f(x)=n-f(x-n)$ alone forces $f\equiv0$, and then $n=1$ gives $0=1-0$, so it has no solutions at all.)

So the error is in the setup, not in the algebra: the ``definition'' of $f$ does not define any function, and from contradictory premises any statement, such as $1=-1$, can be derived. $\boxed{\text{No such function exists.}}$
